\begin{filecontents*}[overwrite]{references.bib}
@article{ramanaharan2025,
  author  = {R. Ramanaharan and D. B. Guruge and J. I. Agbinya},
  title   = {DeepFake video detection: Insights into model generalisation  A Systematic review},
  journal = {Data and Information Management},
  volume  = {9},
  pages   = {100099},
  year    = {2025}
}
@inproceedings{coccomini2022,
  author    = {D. A. Coccomini and R. Caldelli and F. Falchi and C. Gennaro and G. Amato},
  title     = {Cross-forgery analysis of vision transformers and {CNNs} for deepfake image detection},
  booktitle = {Proc. 1st Int. Workshop on Multimedia AI Against Disinformation},
  pages     = {52--58},
  year      = {2022}
}
@inproceedings{lewis2020,
  author    = {J. K. Lewis and others},
  title     = {Deepfake video detection based on spatial, spectral, and temporal inconsistencies using multimodal deep learning},
  booktitle = {IEEE Applied Imagery Pattern Recognition Workshop},
  pages     = {1--9},
  year      = {2020}
}
@article{khormali2022,
  author  = {A. Khormali and J.-S. Yuan},
  title   = {{DFDT}: An end-to-end deepfake detection framework using vision transformer},
  journal = {Applied Sciences},
  volume  = {12},
  number  = {6},
  pages   = {2953},
  year    = {2022}
}
@misc{muppalla2023,
  author = {S. Muppalla and S. Jia and S. Lyu},
  title  = {Integrating audio-visual features for multimodal deepfake detection},
  note   = {arXiv preprint arXiv:2310.03827},
  year   = {2023}
}
@inproceedings{fakeavceleb,
  author    = {H. Khalid and S. Tariq and M. Kim and S. S. Woo},
  title     = {{FakeAVCeleb}: A novel audio-video multimodal deepfake dataset},
  booktitle = {Proc. NeurIPS Datasets and Benchmarks Track},
  year      = {2021}
}
@inproceedings{lavdf,
  author    = {Z. Cai and K. Stefanov and A. Dhall and M. Hayat},
  title     = {Do you really mean that? Content driven audio-visual deepfake dataset and multimodal method for temporal forgery localization},
  booktitle = {Proc. DICTA},
  year      = {2022}
}
@inproceedings{celebdf,
  author    = {Y. Li and X. Yang and P. Sun and H. Qi and S. Lyu},
  title     = {{Celeb-DF}: A large-scale challenging dataset for {DeepFake} forensics},
  booktitle = {Proc. CVPR},
  year      = {2020}
}
@misc{dfdc,
  author = {B. Dolhansky and J. Bitton and B. Pflaum and J. Lu and R. Howes and M. Wang and C. Canton Ferrer},
  title  = {The {DeepFake} Detection Challenge ({DFDC}) dataset},
  note   = {arXiv:2006.07397},
  year   = {2020}
}
@inproceedings{ffpp,
  author    = {A. R{\"o}ssler and D. Cozzolino and L. Verdoliva and C. Riess and J. Thies and M. Nie{\ss}ner},
  title     = {{FaceForensics++}: Learning to detect manipulated facial images},
  booktitle = {Proc. ICCV},
  year      = {2019}
}
@inproceedings{syncnet,
  author    = {J. S. Chung and A. Zisserman},
  title     = {Out of time: Automated lip sync in the wild},
  booktitle = {Proc. ACCV Workshops},
  year      = {2016}
}
@inproceedings{avhubert,
  author    = {B. Shi and W.-N. Hsu and K. Lakhotia and A. Mohamed},
  title     = {Learning audio-visual speech representation by masked multimodal cluster prediction},
  booktitle = {Proc. ICLR},
  year      = {2022}
}
@inproceedings{sbi,
  author    = {K. Shiohara and T. Yamasaki},
  title     = {Detecting deepfakes with self-blended images},
  booktitle = {Proc. CVPR},
  year      = {2022}
}
@inproceedings{univfd,
  author    = {U. Ojha and Y. Li and Y. J. Lee},
  title     = {Towards universal fake image detectors that generalize across generative models},
  booktitle = {Proc. CVPR},
  year      = {2023}
}
@article{shortcut,
  author  = {R. Geirhos and J.-H. Jacobsen and C. Michaelis and R. Zemel and W. Brendel and M. Bethge and F. A. Wichmann},
  title   = {Shortcut learning in deep neural networks},
  journal = {Nature Machine Intelligence},
  volume  = {2},
  pages   = {665--673},
  year    = {2020}
}
@misc{audit,
  title = {What do deepfake benchmarks measure? An audit using frozen self-supervised representations},
  note  = {arXiv:2606.26384},
  year  = {2026}
}
@inproceedings{efficientnet,
  author    = {M. Tan and Q. V. Le},
  title     = {{EfficientNet}: Rethinking model scaling for convolutional neural networks},
  booktitle = {Proc. ICML},
  year      = {2019}
}
@article{mtcnn,
  author  = {K. Zhang and Z. Zhang and Z. Li and Y. Qiao},
  title   = {Joint face detection and alignment using multitask cascaded convolutional networks},
  journal = {IEEE Signal Process. Lett.},
  volume  = {23},
  number  = {10},
  pages   = {1499--1503},
  year    = {2016}
}
@inproceedings{wav2vec2,
  author    = {A. Baevski and Y. Zhou and A. Mohamed and M. Auli},
  title     = {wav2vec 2.0: A framework for self-supervised learning of speech representations},
  booktitle = {Proc. NeurIPS},
  year      = {2020}
}
@misc{dinov3,
  author = {O. Sim{\'e}oni and others},
  title  = {{DINOv3}},
  note   = {arXiv:2508.10104},
  year   = {2025}
}
@misc{mnw,
  author       = {{Microsoft, Northwestern University, and WITNESS}},
  title        = {{MNW} deepfake detection benchmark},
  howpublished = {\url{https://github.com/microsoft/MNW}}
}
@misc{df26,
  title        = {{DF26}},
  howpublished = {Dataset, \url{https://huggingface.co/datasets/DF26/DF26}},
  year         = {2026}
}
@book{bootstrap,
  author    = {B. Efron and R. J. Tibshirani},
  title     = {An Introduction to the Bootstrap},
  publisher = {Chapman \& Hall},
  address   = {New York},
  year      = {1993}
}
@inproceedings{convnext,
  author    = {Z. Liu and H. Mao and C.-Y. Wu and C. Feichtenhofer and T. Darrell and S. Xie},
  title     = {A {ConvNet} for the 2020s},
  booktitle = {Proc. CVPR},
  year      = {2022}
}
@article{adabn,
  author  = {Y. Li and N. Wang and J. Shi and X. Hou and J. Liu},
  title   = {Adaptive batch normalization for practical domain adaptation},
  journal = {Pattern Recognition},
  volume  = {80},
  pages   = {109--117},
  year    = {2018}
}
@misc{pe,
  author = {D. Bolya and others},
  title  = {Perception Encoder: The best visual embeddings are not at the output of the network},
  note   = {arXiv:2504.13181},
  year   = {2025}
}
@inproceedings{clip,
  author    = {A. Radford and others},
  title     = {Learning transferable visual models from natural language supervision},
  booktitle = {Proc. ICML},
  year      = {2021}
}
@inproceedings{hsemotion,
  author    = {A. V. Savchenko},
  title     = {Facial expression and attributes recognition based on multi-task learning of lightweight neural networks},
  booktitle = {Proc. IEEE SISY},
  year      = {2021}
}
\end{filecontents*}

\documentclass[a4paper,12pt,oneside]{report}

\usepackage[
    a4paper,
    left   = 2.54cm,
    right  = 2.00cm,
    top    = 2.00cm,
    bottom = 2.50cm,
    footskip = 1.0cm
]{geometry}

\usepackage{mathptmx}
\usepackage[T1]{fontenc}
\usepackage[utf8]{inputenc}

\usepackage{setspace}
\onehalfspacing

\usepackage{xcolor}
\definecolor{Cardinal}{HTML}{B7202F}
\definecolor{CapeHoney}{HTML}{FEE1A5}
\definecolor{Eclipse}{HTML}{312016}

\usepackage{graphicx}
\usepackage{tabularx}
\usepackage{booktabs}
\usepackage{array}
\usepackage{longtable}
\usepackage{multirow}
\usepackage{amsmath,amssymb}
\usepackage{caption}
\captionsetup{font=small,labelfont={bf,color=Eclipse}}

\usepackage{fancyhdr}
\usepackage{titlesec}
\usepackage[titles]{tocloft}

\usepackage{eso-pic}
\usepackage{tikz}

\usepackage[hidelinks, bookmarksnumbered]{hyperref}

% =============================================================================
% PROJECT METADATA
% =============================================================================
\newcommand{\ProjectTitle}{Multimodal Audio-Visual Deepfake Detection}
\newcommand{\Branch}{Electronics and Computer Engineering}
\newcommand{\BranchShort}{Electronics and Computer Engineering}
\newcommand{\Department}{Department of \BranchShort}
\newcommand{\University}{Somaiya Vidyavihar University}
\newcommand{\College}{K. J. Somaiya School of Engineering}
\newcommand{\Semester}{VII}
\newcommand{\Batch}{2024-25}
\newcommand{\AcademicYear}{2024-2025}
\newcommand{\PlaceCity}{Mumbai-77}

\newcommand{\StudentEntry}[2]{%
    \textbf{\large #1}\\[1pt]
    \normalsize Roll No: #2\\[6pt]
}

\newcommand{\GuideName}{Dr. Tejashree Phatak}
\newcommand{\GuideDesignation}{Associate Professor}

% =============================================================================
% PAGE BORDER
% =============================================================================
\AddToShipoutPictureBG{%
    \begin{tikzpicture}[overlay,remember picture]
        \draw[Cardinal, line width=1.2pt]
            ([xshift= 1.0cm, yshift= 1.0cm]current page.south west)
            rectangle
            ([xshift=-1.0cm, yshift=-1.0cm]current page.north east);
    \end{tikzpicture}%
}

% =============================================================================
% FOOTER
% =============================================================================
\newcommand{\FooterText}{%
    \footnotesize\bfseries
    \Department\ \quad Semester \Semester\ \quad \Batch\ Batch%
}

\pagestyle{fancy}
\fancyhf{}
\renewcommand{\headrulewidth}{0pt}
\renewcommand{\footrulewidth}{1pt}
\fancyfoot[L]{\FooterText}
\fancyfoot[R]{\footnotesize\bfseries Page \thepage}

\makeatletter
\renewcommand{\footrule}{%
    {\color{Cardinal}\hrule \@height\footrulewidth \@width\headwidth}%
    \vskip-\footrulewidth%
}
\makeatother

\fancypagestyle{plain}{%
    \fancyhf{}%
    \renewcommand{\headrulewidth}{0pt}%
    \renewcommand{\footrulewidth}{1pt}%
    \fancyfoot[L]{\FooterText}%
    \fancyfoot[R]{\footnotesize\bfseries Page \thepage}%
}

\fancypagestyle{titlepg}{%
    \fancyhf{}%
    \renewcommand{\headrulewidth}{0pt}%
    \renewcommand{\footrulewidth}{0pt}%
}

% =============================================================================
% HEADINGS
% =============================================================================
\titleformat{\chapter}[display]
    {\normalfont\Huge\bfseries\color{Cardinal}}
    {\chaptertitlename\ \thechapter}{14pt}{\Huge}
\titlespacing*{\chapter}{0pt}{-20pt}{24pt}

\titleformat{\section}
    {\normalfont\Large\bfseries\color{Eclipse}}{\thesection}{1em}{}
\titleformat{\subsection}
    {\normalfont\large\bfseries\color{Eclipse}}{\thesubsection}{1em}{}
\titleformat{\subsubsection}
    {\normalfont\normalsize\bfseries\color{Eclipse}}{\thesubsubsection}{1em}{}

\renewcommand{\cfttoctitlefont}{\Huge\bfseries\color{Cardinal}}
\renewcommand{\cftloftitlefont}{\Huge\bfseries\color{Cardinal}}
\renewcommand{\cftlottitlefont}{\Huge\bfseries\color{Cardinal}}
\renewcommand{\cftchapfont}{\bfseries}
\renewcommand{\cftchappagefont}{\bfseries}
\setlength{\cftbeforechapskip}{6pt}

\setlength{\parskip}{6pt}
\setlength{\parindent}{0pt}

% =============================================================================
% BEGIN DOCUMENT
% =============================================================================
\begin{document}

% --
% TITLE PAGE
% --
\pagenumbering{roman}
\setcounter{page}{1}
\thispagestyle{titlepg}

\begin{singlespace}
\begin{center}

    \includegraphics[width=4.8cm]{1777459645157_kjsce.png}

    \vspace{8pt}
    {\Large\bfseries\color{Eclipse}\University\par}
    \vspace{2pt}
    {\large\bfseries \College\par}
    \vspace{2pt}
    {\large\bfseries \Department\par}

    \vspace{22pt}
    {\huge\bfseries\color{Cardinal} \ProjectTitle\par}

    \vspace{18pt}
    {\normalsize Submitted In Partial Fulfillment of Requirements\par}
    {\normalsize For the Degree Of\par}
    \vspace{6pt}
    {\Large\bfseries Bachelor of Technology\par}
    {\large\bfseries (\Branch)\par}

    \vspace{16pt}
    {\normalsize By\par}
    \vspace{6pt}

    \StudentEntry{Vaibhav Shiroorkar}{16014024805}
    \StudentEntry{Adithya Jithesh}{16014024803}
    \StudentEntry{Shivansh Gupta}{16014023052}


    \vspace{4pt}
    {\normalsize Guide\par}
    \vspace{2pt}
    {\Large\bfseries \GuideName \par}
    {\normalsize \GuideDesignation\par}

    \vspace{18pt}
    {\large\bfseries \University\par}
    {\normalsize Vidyavihar, Mumbai - 400 077\par}
    \vspace{4pt}
    {\large\bfseries 2026-2027\par}
\end{center}
\end{singlespace}

\clearpage

% --
% CERTIFICATE
% --
\setcounter{page}{1}

\begin{center}
    {\Large\bfseries\color{Eclipse} \University\par}
    \vspace{4pt}
    {\Large\bfseries \College\par}
    \vspace{30pt}
    {\Huge\bfseries\color{Cardinal} Certificate\par}
\end{center}

\vspace{16pt}

This is to certify that the dissertation report entitled
\textbf{``\ProjectTitle''} submitted by the following students at the end of
Semester \Semester\ of LY B.\,Tech is a bona fide record for partial
fulfillment of requirements for the degree of \textbf{Bachelor of Technology
(\Branch)} of \University.

\vspace{12pt}

\begin{tabular}{@{}p{0.5\textwidth}p{0.45\textwidth}@{}}
    \textbf{Vaibhav Shiroorkar} \quad (16014024805) & \\[8pt]
    \textbf{Adithya Jithesh} \quad (16014024803) & \\[8pt]
    \textbf{Shivansh Gupta} \quad (16014023052) & \\
\end{tabular}

\vspace{60pt}

\noindent
\begin{tabular}{@{}p{0.45\textwidth}@{\hspace{0.05\textwidth}}p{0.45\textwidth}@{}}
    \rule{6cm}{0.4pt} & \rule{6cm}{0.4pt} \\
    \textbf{Guide} & \textbf{Head of the Department} \\[2pt]
    \GuideName & Dr. Jagannath Nirmal \\[20pt]

    \rule{6cm}{0.4pt} & \rule{6cm}{0.4pt} \\
    \textbf{Examiner} & \textbf{Director} \\[2pt]
    & \\
\end{tabular}

\vspace{30pt}

\noindent\textbf{Date:}\\
\noindent\textbf{Place:} \PlaceCity

\clearpage

% --
% DECLARATION
% --
\begin{center}
    {\Large\bfseries\color{Eclipse} \University\par}
    \vspace{4pt}
    {\Large\bfseries \College\par}
    \vspace{30pt}
    {\Huge\bfseries\color{Cardinal} Declaration\par}
\end{center}

\vspace{12pt}

We declare that this written report submission represents the work done based
on our and / or others' ideas with adequately cited and referenced the
original source. We also declare that we have adhered to all principles of
intellectual property, academic honesty and integrity as we have not
misinterpreted or fabricated or falsified any idea / data / fact / source /
original work / matter in our submission.

\vspace{6pt}

We understand that any violation of the above will be cause for disciplinary
action by the college and may evoke the penal action from the sources which
have not been properly cited or from whom proper permission is not sought.

\vspace{30pt}

\renewcommand{\arraystretch}{2.4}
\begin{tabular}{|p{0.45\textwidth}|p{0.45\textwidth}|}
    \hline
    \textbf{Signature of the Student} \newline\rule{6cm}{0.4pt}\newline
    \textbf{Roll No.} \rule{4cm}{0.4pt} &
    \textbf{Signature of the Student} \newline\rule{6cm}{0.4pt}\newline
    \textbf{Roll No.} \rule{4cm}{0.4pt} \\ \hline
    \textbf{Signature of the Student} \newline\rule{6cm}{0.4pt}\newline
    \textbf{Roll No.} \rule{4cm}{0.4pt} &
    \textbf{Signature of the Student} \newline\rule{6cm}{0.4pt}\newline
    \textbf{Roll No.} \rule{4cm}{0.4pt} \\ \hline
\end{tabular}
\renewcommand{\arraystretch}{1.0}

\vspace{24pt}

\noindent\textbf{Date:}\\
\noindent\textbf{Place:} \PlaceCity

\clearpage

% --
% ABSTRACT
% --
\chapter*{Abstract}
\addcontentsline{toc}{chapter}{Abstract}

The question behind this report is a simple one: is the model we have now
actually better than the one we started with, and if so, better in what way?
The project began on what is now the \texttt{old} branch of our repository as a
five-stream audio-visual design. It had three visual backbones (Xception,
EfficientNet-B0 and a self-supervised transformer that was supposed to be
DINOv3), a lip-sync stream built on AV-HuBERT and Whisper, an
emotion-consistency stream built on HSEmotion and Wav2Vec2, and a
feature-level fusion layer on top of all five. When we went through that
branch's code rather than its planning documents, we found that only the three
visual streams had ever been trained and tested. They scored 0.996 ROC AUC
(EfficientNet-B0), 0.990 (Xception) and a suspicious 1.000 for the
transformer, whose checkpoint on disk is actually labelled DINOv2, not DINOv3.
The lip-sync stream, the emotion stream and the fusion layer were empty
folders. Nothing checked where a checkpoint came from, and no stream had ever
been tested on a second dataset.

In the current phase, on \texttt{main}, we first cut the scope down to a single
visual branch, EfficientNet-B0 with a GRU, and put it through checks the old
branch never had. These were four extra training runs to rule out a lucky
seed, a cross-dataset test on Celeb-DF-v2, and a six-value provenance check
(checkpoint hash, MLflow run ID, split hash, training commit, seed and
preprocessing hash) that has to pass before a checkpoint can give any verdict.
Its in-distribution score of 0.999 fell to 0.6476 on Celeb-DF-v2. With that
baseline in place we went back and built the multimodal system the old branch
had planned, and then tested whether each part does what its name says. For
the most part it does not. Fusion adds $+0.0159$ ROC AUC in-domain, but on
DFDC, a corpus it never saw, it is no better than the visual branch at any of
three seeds and is measurably worse at two. A cross-attention emotion stream
reached 0.9991 ROC AUC while its attention never moved away from the uniform
value. The visual models had partly learned the video codec: re-encoding
genuine clips with the codec of the largest fake category makes 93.3 percent
of them look fake. A checkpoint chosen on validation loss scored below chance,
and the training seed alone shifted cross-corpus ROC AUC by 0.1078. We caught
each of these by logging a quantity next to the training objective without
ever optimizing it. So the improvement documented here is less about accuracy
and more about knowledge: the current project knows how good its models really
are, and where they break.

\clearpage

% --
% ACKNOWLEDGEMENTS
% --
\chapter*{Acknowledgements}
\addcontentsline{toc}{chapter}{Acknowledgements}

We would like to thank our guide, \textbf{\GuideName}, for her help throughout
this project. That help went well beyond formal reviews. When we were stuck she
pointed us in a useful direction, and when it was better for us to work
something out ourselves she gave us room to do that. Looking back at our own
earlier work and admitting where it fell short was not always comfortable, and
her encouragement to write openly about what worked and what did not made this
a better report than a purely positive one would have been.

We are also grateful to the faculty and staff of the \Department\ at \College\
for the resources and support behind this work, including access to the GPU
hardware that both phases of the project relied on. Finally, thanks to our
families and friends for keeping us sane through the long waits on training
runs, the CUDA errors, and the evenings spent reading through our own old
commits trying to remember what we had actually built.

\clearpage

% --
% TABLE OF CONTENTS
% --
\renewcommand{\contentsname}{Contents}
\tableofcontents
\clearpage

\listoftables
\addcontentsline{toc}{chapter}{List of Tables}
\clearpage

\listoffigures
\addcontentsline{toc}{chapter}{List of Figures}
\clearpage

% --
% NOMENCLATURE
% --
\chapter*{Nomenclature}
\addcontentsline{toc}{chapter}{Nomenclature}

\begin{tabularx}{\textwidth}{@{}>{\raggedright\arraybackslash}p{3.8cm}X@{}}
    \toprule
    \textbf{Term} & \textbf{Meaning} \\
    \midrule
    Legacy phase       & The project as it existed on the \texttt{old} branch of the repository \\
    Current phase      & The project as it exists on the \texttt{main} branch of the repository \\
    Design A           & The current phase's three-branch system (visual, audio, sync) combined by late, calibrated logistic-regression fusion \\
    Design B           & The current phase's configurable-stream system combined by a feature-level fusion head \\
    Stream             & One backbone or signal-processing path that produces a clip-level embedding \\
    Cross-modal stream & A stream that compares audio against video directly, via cross-attention \\
    Fusion             & The layer that combines every stream's embedding into one final verdict \\
    Diagonal attention mass & The share of a cross-attention map that falls near its diagonal; used to check whether a stream actually aligns sound with mouth movement \\
    ROC AUC            & Area under the receiver operating characteristic curve \\
    EER                & Equal error rate, the point where false positives equal false negatives \\
    Identity-disjoint split & A train/val/test split in which no single person appears in more than one partition \\
    In-distribution    & Evaluated on a held-out slice of the same dataset used for training \\
    Cross-dataset / generalization & Evaluated on a dataset never touched during training or tuning \\
    DFDC               & The DeepFake Detection Challenge dataset, the main cross-corpus test set in Chapter~\ref{ch:multimodal} \\
    Codec shortcut     & A model separating real from fake by the video encoder used rather than by the manipulation itself \\
    Paired bootstrap interval & A confidence interval for the difference between two systems, computed on the same resampled identities for both \\
    Out-of-fold        & Predictions made by a model on data it was not trained on, used so that a fusion head is never fitted on leaked scores \\
    Seed spread        & The range of a metric across training runs that differ only in their random seed \\
    Provenance check   & Automatic verification that a checkpoint matches its recorded hash, run ID, split, commit, and seed \\
    Saturated prediction & A model output so close to 0 or 1 that it behaves like a hard label rather than a calibrated probability \\
    Software fixture   & A placeholder artifact built to satisfy an interface, not a trained research model \\
    W\&B               & Weights and Biases, the experiment tracker used on the legacy phase \\
    MLflow             & The experiment tracker used on the current phase \\
    \bottomrule
\end{tabularx}

\clearpage

% =============================================================================
% MAIN MATTER
% =============================================================================
\pagenumbering{arabic}
\setcounter{page}{1}

% --
\chapter{Introduction}
\label{ch:introduction}

\section{Why This Report Exists}

Any project that runs for a while builds up a history, and that history tends
to get squeezed into one line such as ``we improved the model.'' In our case
that line is true, but it leaves out the more interesting part. The project
did not just get better at the same job; it changed what it was trying to
show. The \texttt{old} branch was an attempt to build the most capable detector
we could, as quickly as we could, using as many kinds of signal as possible.
The work on \texttt{main} is an attempt to build detectors whose reported
performance can be trusted, even when that performance turns out lower than we
hoped. We think that change of goal is the real result of the project, and this
report is organised around it.

\section{How This Report Was Built}

Chapters \ref{ch:legacy}, \ref{ch:current} and \ref{ch:multimodal} are written
from each branch's actual code, checkpoints, logs and documentation, not from
memory or from summaries written later. The \texttt{old} branch's evaluation
numbers come from \texttt{logs/test\_eval\_v4.txt}, a held-out test run dated
2026-08-04. The numbers for \texttt{main} come from its training history, its
validation metrics files and its MLflow store, checked against the checkpoint
hashes and run IDs recorded there. The multimodal results in
Chapter~\ref{ch:multimodal} are registered results, each carrying the SHA-256
hashes of the artifacts that produced it. Wherever a document on either branch
said something the code did not back up, we point it out. One example: the
\texttt{old} branch's architecture document calls one stream ``DINOv3'', while
its checkpoint folder and training logs call it ``dinov2''. We report the
mismatch instead of choosing whichever name sounds better.

\section{Scope}

The visual detection work on both branches gets the most detail, since on the
old branch it is the only part that was ever trained and evaluated. The old
branch's lip-sync, emotion and fusion components are described as what they
were: planned, but never built. For the current phase we cover the
cross-dataset evaluation and provenance system in full, because that is the
capability that did not exist before. We also cover the later multimodal
experiments, in which the current phase built the audio-visual streams and
fusion head the old branch only planned, and then tested whether they actually
help.

\section{How to Read This Report}

Chapter \ref{ch:literature} places the work within the wider deepfake
detection literature, so the reader has some sense of what a trustworthy
result looks like before either branch comes up. Chapter \ref{ch:legacy}
describes the legacy phase as its code shows it, not as its planning documents
describe it. Chapter \ref{ch:current} does the same for the current phase's
visual baseline. Chapter \ref{ch:multimodal} covers the multimodal experiments
built on top of that baseline and the shortcuts and validation traps they
exposed. Chapter \ref{ch:comparison} puts the phases side by side and answers
the main question directly. Chapter \ref{ch:verification} separates what is
verified from what is still blocked or pending. Chapter \ref{ch:future} lists
the next steps, and Chapter \ref{ch:conclusions} closes the report.

% --
\chapter{Literature Review}
\label{ch:literature}

\section{Introduction}

Before comparing two phases of our own project, it is worth setting an outside
reference point. What does the deepfake detection community already know about
the gap between a model's reported accuracy and how well it holds up in
practice? We keep this review short. It is not meant as a full survey; we
include it because almost every finding below is something our project ended up
rediscovering on its own.

\section{Datasets and the Standard Transfer Test}

Most audio-visual deepfake work relies on a handful of public corpora.
FakeAVCeleb~\cite{fakeavceleb} and LAV-DF~\cite{lavdf} pair manipulated faces
with manipulated voices, which makes them the natural starting point for any
detector that reads both modalities. The usual way to test transfer is to train
on FaceForensics++~\cite{ffpp} and then evaluate on other corpora such as
Celeb-DF~\cite{celebdf} and DFDC~\cite{dfdc}. Both phases of our project train
mainly on FakeAVCeleb, and the current phase uses Celeb-DF-v2, DFDC, LAV-DF and
FaceForensics++ in the roles described in Chapters~\ref{ch:current}
and~\ref{ch:multimodal}.

\section{The Generalization Gap Is a Known Problem}

Ramanaharan et al. reviewed 108 published deepfake detection studies and found
that only 46.3\% of them considered their own model generalisable
\cite{ramanaharan2025}. Before our project had any cross-dataset number of its
own, this read as a useful but fairly abstract warning. After watching our own
current-phase checkpoint fall from 0.999 to 0.6476 ROC AUC just by switching
the evaluation dataset, it read more like a prediction. It also describes our
legacy phase quite well: three trained streams, each tested only on data from
its training distribution, is exactly the pattern the review found across the
field.

\section{Cross-Forgery and Cross-Dataset Sensitivity}

Coccomini et al. trained detectors on one manipulation method and tested them
on another. Both convolutional networks and Vision Transformers lost a
noticeable amount of accuracy under that shift \cite{coccomini2022}.
FakeAVCeleb and Celeb-DF-v2, the two datasets in our first cross-dataset test,
differ in exactly the ways such studies flag as important: different subjects,
different manipulation pipelines and different capture and compression
conditions. The drop we measured, roughly 35 points of ROC AUC, falls well
within what this body of work would predict for a mismatch of that kind. That
is one reason we treat the drop as a real finding and not something to explain
away.

\section{Temporal and Multimodal Signals}

Lewis et al. combined spatial, spectral and temporal signals and reported
accuracy of up to 98\% on controlled benchmarks \cite{lewis2020}. Muppalla et
al. showed that audio-visual fusion can catch manipulations that a visual-only
model misses, since faking two channels convincingly at the same time is
harder than faking one \cite{muppalla2023}. Khormali and Yuan's DFDT work
argues that transformer-based temporal modelling helps generalization, not just
in-distribution accuracy \cite{khormali2022}. Mouth-audio synchronization in
particular has been modelled since SyncNet~\cite{syncnet}, and
AV-HuBERT~\cite{avhubert} learns a joint audio-visual speech representation
from unlabelled video. Taken together, these papers are close to a restatement
of the reasoning behind our original five-stream design: combine several visual
backbones with real cross-modal signal, because no single stream is likely to
catch every manipulation family by itself. The literature supports that idea.
What it does not support is trusting any single stream's number before it has
been through seed repetition and cross-dataset testing. Our legacy phase
skipped that step, and judging by Ramanaharan et al.'s review, so did a fair
share of published work.

\section{Training Choices Aimed at Transfer}

Two other lines of work shaped the experiments in Chapter~\ref{ch:multimodal}.
Self-blended images~\cite{sbi} train a face forgery detector on blending
artifacts instead of on the output of any one generator, with the aim of
transferring to generators it has never seen. Separately, Ojha et al. showed
that a frozen vision-language encoder with only a linear probe on top transfers
well across image generators~\cite{univfd}. We use the first idea as a training
recipe for our visual branch, and the second as the basis for our experiments
on fully generated video.

\section{Shortcut Learning}

Geirhos et al. describe shortcut learning as a model solving a benchmark
through a cue the task designers never intended~\cite{shortcut}. Most of the
checks in Chapter~\ref{ch:multimodal} are really searches for shortcuts of this
kind. Examples include a detector reading the video codec instead of the
manipulation, or a classifier separating real from generated clips by whether
an audio track is present. A recent audit of deepfake benchmarks using frozen
self-supervised features~\cite{audit} covers ground close to our own finding
about choosing the right region of the frame to look at.

\section{Where Our Project Fits}

Put simply, the legacy phase is not unusual for the field. It is a small, local
example of a pattern the literature describes as common: real architectural
ambition and real trained components, but no cross-dataset check on any of
them. The current phase narrowed the scope, added seed repetition and a real
cross-dataset test, and then rebuilt the multimodal design under those same
checks. That is our attempt to end up on the right side of Ramanaharan et al.'s
46.3\% figure. The next three chapters describe how that attempt went.

% --
\chapter{The Legacy Architecture}
\label{ch:legacy}

\section{The Plan}

The \texttt{old} branch's main planning document,
\texttt{docs/PROJECT\_OVERVIEW.md}, sets out an ambitious five-stream design.
Three visual-only streams (Xception, EfficientNet-B0 and a self-supervised
vision transformer) would each read a sequence of face crops and produce a
clip-level embedding without any access to the audio. Two cross-modal streams
would compare audio and video directly. One was a lip-sync stream using
AV-HuBERT for video and Whisper for audio. The other was an
emotion-consistency stream using HSEmotion for facial expression and Wav2Vec2
for vocal affect. Both would use cross-attention so that one modality could
query the other. All five streams would write to a shared feature store, and a
fusion MLP over the concatenated embeddings, rather than a simple average of
five scores, would produce the final verdict.

It was a well-reasoned design, and the document explaining it handles real
trade-offs with care. For example, it treats frozen backbones versus
end-to-end fine-tuning as an open question to settle against the team's
actual GPU budget. It also points out, correctly, that a face swap and a
lip-sync repaint leave evidence in different places, which is exactly why a
detector relying on one signal struggles against a varied set of attacks.
Figure~\ref{fig:arch-overview} shows the full planned design.

\begin{figure}[htbp]
    \centering
    \includegraphics[width=0.60\textwidth,keepaspectratio]{fig_arch_overview.png}
    \caption{The legacy phase's planned five-stream audio-visual architecture: an input clip is split into visual and audio pre-processing, feeding three visual streams (the self-supervised transformer and EfficientNet, with a temporal sequence model) and two cross-modal streams (lip-sync via AV-HuBERT + Wav2Vec, emotion via HSEmotions + Wav2Vec2, each through a cross-attention mechanism), all combined by a fusion layer before the final decision and explainability stage.}
    \label{fig:arch-overview}
\end{figure}

\section{What Was Actually Built}

The plan and the code turn out to tell different stories. Where they disagree,
we go with the code.

\subsection{Three Visual Streams, Genuinely Trained}

Checkpoint files exist on disk for three visual backbones, and a held-out test
evaluation dated 2026-08-04 scored all three under the same conditions.
Table~\ref{tab:legacy-results} reproduces that evaluation as it was recorded.
Continuing from the pre-processing flow in Figure~\ref{fig:arch-overview},
Figure~\ref{fig:stage2a} shows the Stage 2A visual-stream design, in which two
backbones (EfficientNet and DINOv3) each feed a BiLSTM temporal head. Xception,
the third stream that was eventually trained, does not appear in the diagram.

\begin{figure}[htbp]
    \centering
    \includegraphics[width=0.60\textwidth,keepaspectratio]{fig_stage2a_visual_new.png}
    \caption{Stage 2A visual-stream design: parallel EfficientNet and DINOv3 backbones, each producing an embedding fed to a BiLSTM temporal sequence model before the shared fusion layer.}
    \label{fig:stage2a}
\end{figure}

These are real numbers from a real evaluation, and it would be unfair to the
earlier work to pretend otherwise. Still, the evaluation log itself raises
doubts about them, especially the third row. It states that the perfect score
is ``an in-distribution ceiling, not evidence of generalization.'' It also notes
that all three streams were heavily saturated: most predictions sat above 0.99
or below 0.01, with almost nothing in the uncertain middle. That makes the
fixed 0.5 threshold close to meaningless as an operating point. The same log
recommends fitting a temperature-scaling parameter for each stream before
building fusion. That step was never done, and no later run in the repository
comes back to it.

\begin{table}[ht]
    \centering
    \caption{Legacy branch: held-out test results for all three trained visual streams, 300 clips (80 real, 220 fake), identity-disjoint from train and validation.}
    \label{tab:legacy-results}
    \begin{tabularx}{\textwidth}{Xrrrr}
        \toprule
        \textbf{Stream} & \textbf{Accuracy} & \textbf{ROC AUC} & \textbf{EER} & \textbf{Log loss} \\
        \midrule
        EfficientNet-B0 & 0.977 & 0.996 & 0.018 & 0.107 \\
        Xception        & 0.943 & 0.990 & 0.032 & 0.178 \\
        Self-supervised transformer (checkpoint labeled \texttt{dinov2}) & 1.000 & 1.000 & 0.000 & 0.0001 \\
        \bottomrule
    \end{tabularx}
\end{table}

\subsection{A Naming Problem Worth Stating Plainly}

The planning document calls the third visual stream ``DINOv3'', and its code
lives in a folder named \texttt{models/streams/dinov3/}. The checkpoint that
was actually trained and evaluated, however, is
\texttt{checkpoints/dinov2/best.pt}, and both the training and evaluation logs
call it \texttt{dinov2} throughout. This is more than a typo. It means the
legacy documentation was already describing a slightly different model from the
one that was run, months before this report was written. Chapter
\ref{ch:current} argues that the current branch is a real improvement partly
because its provenance system is designed to stop exactly this kind of drift
between what a document says and what a checkpoint is.

\subsection{Two Streams and a Fusion Layer That Were Never Built}
\label{sec:legacy-unbuilt}

Both cross-modal streams and the fusion layer exist in the repository only as
folders containing a single empty \texttt{\_\_init\_\_.py} file. We confirmed
this by inspecting the folders directly, not by reading the documentation.
The folders \texttt{models/streams/lipsync/}, \texttt{models/streams/emotion/}
and \texttt{fusion/} are all empty. There is no AV-HuBERT,
Whisper, HSEmotion or emotion-side Wav2Vec2 integration, no cross-attention
mechanism and no fusion MLP anywhere in the legacy codebase, at any commit we
examined. The branch's dashboard is upfront about this. Its changelog
describes the Fusion and Explainability pages as ``locked, listed with a lock
icon, no interactive controls,'' which matches the code. Figures~\ref{fig:stage2b}
and~\ref{fig:stage3} show these components as they were designed but never
coded on this branch.

\begin{figure}[htbp]
    \centering
    \includegraphics[width=0.75\textwidth,keepaspectratio]{fig_stage2b_crossmodal.png}
    \caption{Stage 2B cross-modal streams as planned: a lip-sync stream (AV-HuBERT video encoder, Wav2Vec2 audio encoder) and an emotion stream (HSEmotion video encoder, Wav2Vec2 audio encoder), each fusing modalities through temporal cross-attention into fused vectors before the fusion layer. Both streams were planned but, as this section records, never implemented beyond an empty directory.}
    \label{fig:stage2b}
\end{figure}

\begin{figure}[htbp]
    \centering
    \includegraphics[width=0.58\textwidth,keepaspectratio]{fig_stage3_fusion.png}
    \caption{Stage 3 fusion and decision as planned: every stream's embedding enters a shared feature store, is concatenated, passed through a multi-layer perceptron and a sigmoid activation to a final score between 0 and 1, and turned into a real-or-fake decision followed by an explainability stage. This fusion layer was never built.}
    \label{fig:stage3}
\end{figure}

\section{Methodological Gaps}

Apart from the unfinished streams, three gaps in the legacy branch's methods
matter for the comparison in Chapter \ref{ch:comparison}.

\textbf{No seed repetition.} Each of the three visual streams was trained only
once. Nothing in the legacy branch shows any stream being retrained with a
different random seed to see whether its score reflected the architecture or a
lucky draw. The perfect transformer score in Table \ref{tab:legacy-results} is
the result a second run would have helped most with, and that run never
happened.

\textbf{No cross-dataset evaluation.} Every number in
Table~\ref{tab:legacy-results} was measured on a held-out split of the same
dataset every stream was trained on, FakeAVCeleb. \texttt{PROJECT\_OVERVIEW.md} does
name Deepfake-Eval-2024 as a dataset ``held aside for the final real-world
generalization test, never used for training,'' so the intention was right.
But there is no evaluation against it, or any other dataset, anywhere in the
branch's logs, checkpoints or scripts. The plan to measure generalization was
written down; the measurement was never made.

\textbf{No provenance checking.} Checkpoints are loaded by file path alone.
There is no hash check and no record linking a checkpoint file to the commit,
data split or seed that produced it. The branch's notes do mention a narrower
issue: checkpoints have to be saved as plain dictionaries, because the
dashboard's \texttt{weights\_only=True} loading will not unpickle anything else.
That is a serialization rule, not a provenance system. Nothing on the branch
would have caught the DINOv2/DINOv3 mix-up, because nothing ever checked a
loaded checkpoint against what the documents said about it.

\section{A Note on Process}

The way the legacy branch was built also deserves a mention, because the
pattern tells us something. \texttt{PROJECT\_OVERVIEW.md} records a
``deliberate reset'' on 2026-07-23, in which the whole Stage 1 and Stage 2
implementation was deleted and rebuilt from scratch. It also records a
contradiction over class balance in validation and test. The original plan was
to keep the dataset's natural ratio of roughly 40 fakes to every real clip in
those splits. A later decision reversed that and undersampled every split,
including validation and test, to a fixed 1-to-3 ratio. The document marks this
reversal with a warning symbol instead of hiding it.

None of this is a criticism of the people doing the work. Building a
five-stream multimodal system from nothing, in a small team new to deep
learning, with a 6 GB laptop GPU as the limiting factor, is a hard thing to
pull off. If anything, a planning document that openly records its own resets
and contradictions shows a team being straight with itself. The current
branch's provenance system takes that same habit of catching and writing down
inconsistencies and makes it automatic, so it no longer depends on someone
remembering to update a document.

\section{Dataset and Splitting Strategy}

On its own terms, the legacy phase's data strategy was carefully thought
through. FakeAVCeleb v1.2~\cite{fakeavceleb} was picked as the main dataset
because it contains real audio manipulation alongside visual manipulation,
across four categories: RealVideo-RealAudio, RealVideo-FakeAudio,
FakeVideo-RealAudio and FakeVideo-FakeAudio. Its manipulation-method column
(real, rtvc, faceswap, fsgan, wav2lip and two layered combinations) also lines
up neatly with the three detection principles behind the five-stream design.
Splits were made at the identity level. Without that, a real clip and its
fake version, which share background, lighting and framing, could land on
opposite sides of a split, letting the model exploit the overlap instead of
learning real manipulation cues. The current phase follows the same
identity-disjoint approach; it is one of the few methodological choices that
carried over unchanged.

The validation and test class ratio is where things became contested.
\texttt{PROJECT\_OVERVIEW.md} lists this as an unresolved issue. The stated
intent was for validation and test to keep the dataset's natural, heavily
imbalanced distribution of roughly 40 fakes per real clip. The splitting
script that actually shipped, however, undersampled fakes to 1-to-3 in every
split, so that precision, recall and F1 would stay stable on a 300-clip
evaluation set. The results in Table~\ref{tab:legacy-results} were measured
under that 1-to-3 ratio. Anyone comparing them with a different evaluation
protocol, including the current phase's validation numbers, needs to keep this
in mind. It is also why Chapter~\ref{ch:comparison} does not treat the two
branches' raw accuracy figures as directly comparable.

\section{Preprocessing}

The legacy preprocessing pipeline was rebuilt from a snapshot commit as a set
of shared, pure functions under \texttt{preprocessing/ops/}. It produced
224$\times$224 face crops from a margin-padded bounding box and 96$\times$96
mouth crops located from facial landmarks, and stored a per-frame detection
confidence alongside each crop. Audio was extracted as mono at 16 kHz.

FakeAVCeleb has a known quirk: fake-audio clips often start with extra
silence, which a model could learn to use as a shortcut instead of a real
manipulation cue. The pipeline measured this silence for each clip and
offset the start of frame and audio sampling to cancel it, keeping the two
modalities aligned. A five-point face alignment step, which pose-normalised
every crop onto a standard template, was added during this phase and described
in the project's notes as the ``main SOTA gain''. It was later removed again,
and \texttt{docs/ideas.md} records the reversal. Adding a technique, keeping it
long enough to measure, and dropping it when the effect does not hold up is a
good research habit, and it deserves credit even though the cross-modal streams
meant to use this pipeline were never built.

\section{Tooling Decisions}

Three tools were given three separate jobs, and the project documentation is
clear that this division was settled and not up for casual debate. Streamlit
had two uses only. One was a preprocessing dashboard for adjusting settings
such as silence trimming, frame sampling rate and face-crop margin on a sample
of ten to twenty clips. The other was a read-only viewer that pulled finished
runs from Weights and Biases. Streamlit never ran or contained a training loop.
An earlier version of the dashboard did have a tab that generated a
background-trainer command, but it was deleted because, in the project's own
words, ``a command builder is not a training feature,'' and its presence
suggested that training lived in the dashboard when it did not.

Weights and Biases tracked every actual training run. Whenever backbone choice,
backbone freezing, learning rate, batch size or preprocessing parameters were
being compared, those settings had to be logged in the run configuration.
Per-category validation accuracy was also logged every epoch, so that a single
aggregate number could not hide a whole manipulation family being missed.
Training ran as a background script or a W\&B sweep agent, never as a notebook
cell that would die with a dropped connection.

\section{Compute Environment}

The legacy phase's compute assumptions explain several of its architecture
choices, so they are worth recording. The main development machine had an
NVIDIA RTX 5070 Ti with 15.9 GB of VRAM, verified as of 2026-07-23. The
project notes make clear, though, that a teammate's RTX 3060 laptop GPU with
only 6 GB of VRAM was treated as the binding constraint for batch size and
architecture decisions. That is why training used a batch size of 2 with
gradient accumulation to an effective batch of 16, why face detection in the
pre-caching workers ran on the CPU, and why EfficientNet-B0, the lightest of
the three planned visual backbones, was built first. The documentation also
admits that an earlier draft claimed there was ``no compute limit'', and a
later one treated 6 GB as the limit for everyone. It says both were wrong: the
first was wishful, and the second generalised from one person's laptop without
checking whether it applied across the team. Corrections like this, sitting in
the project's own version history, are part of why we see the legacy phase's
shortcomings as gaps in a difficult, ambitious project and not as carelessness.

\section{Work Order and Stage Plan}

The legacy documentation sets out a ten-stage build order, shown in
Table~\ref{tab:stageplan}. It starts with the shared data pipeline, moves
through each stream in turn, and ends with fusion, evaluation and
explainability. Comparing the table with the findings above shows how far the
work got. Stage 1 and the EfficientNet-B0 part of Stage 2 were completed and
verified. Stage 3 produced two more trained streams (Xception and the
transformer labelled \texttt{dinov2}). Stages 4 to 10 never went beyond the
empty folders described in Section~\ref{sec:legacy-unbuilt}.

\begin{table}[ht]
    \centering
    \caption{Legacy phase's planned ten-stage build order, against what was actually reached.}
    \label{tab:stageplan}
    \begin{tabular}{clp{4.6cm}}
        \toprule
        \textbf{Stage} & \textbf{Content} & \textbf{Status} \\
        \midrule
        1  & Data pipeline & Complete, rebuilt after the 2026-07-23 reset \\
        2  & First visual stream (EfficientNet-B0) & Complete, trained and evaluated \\
        3  & Remaining visual streams (Xception, transformer) & Complete, trained and evaluated \\
        4  & Lip-sync cross-modal stream & Not started, empty directory only \\
        5  & Emotion cross-modal stream & Not started, empty directory only \\
        6  & Fusion & Not started, empty directory only \\
        7  & Ablation support & Not started \\
        8  & Self-supervised pretraining (stretch) & Not started \\
        9  & Full evaluation & Not started; no cross-dataset run exists \\
        10 & Explainability and write-up & Not started; dashboard shows this locked \\
        \bottomrule
    \end{tabular}
\end{table}

% --
\chapter{The Current Architecture}
\label{ch:current}

\section{The Decision to Narrow Scope}

The biggest change between the two branches was a decision about scope, not a
technical change. Instead of continuing to build outward across five streams,
we took the one stream that had already been trained, evaluated and
checkpointed, the visual branch, and pushed it through a much harder set of
questions. Is its score stable when it is retrained? Does it hold up on a
dataset it has never seen? Answering both takes real extra work, and neither
answer makes the headline accuracy look better. That is probably why neither
was asked on the legacy branch, and why answering them is the current
branch's main contribution. Only after this baseline was in place did we go
back to the multimodal plan, as described in Chapter~\ref{ch:multimodal}.

\section{Architecture}

The current visual model combines an EfficientNet-B0
backbone~\cite{efficientnet} with a GRU temporal head. For each clip it reads
a sequence of 16 face crops, each $224\times224$ pixels, detected and tracked
with MTCNN~\cite{mtcnn}, instead of judging a single frame on its own. It is
close to the legacy branch's EfficientNet-B0 stream, which was the strongest of
the three properly evaluated legacy streams. The main difference is the
temporal head: the legacy BiLSTM was swapped for a GRU, which is lighter and
let us iterate faster across the retraining runs in
Section~\ref{sec:seed-repetition}.

\section{Ruling Out a Lucky Training Run}
\label{sec:seed-repetition}

Before trusting the frozen checkpoint's validation score, we trained four more
runs on the same cached data, changing only the random seed, the learning rate
or the number of epochs. Table~\ref{tab:seed-comparison} shows the results.
This directly addresses the gap noted in Chapter~\ref{ch:legacy}: the legacy
branch never repeated a training run, and its most striking result, the
transformer's perfect 1.000 AUC, was never checked against this kind of
variation.

\begin{table}[ht]
    \centering
    \caption{Current branch: repeating the training recipe across seeds and hyperparameters, all evaluated on the same 400 source-disjoint FakeAVCeleb validation rows.}
    \label{tab:seed-comparison}
    \begin{tabular}{lrrrr}
        \toprule
        \textbf{Run} & \textbf{Seed} & \textbf{Val. loss} & \textbf{ROC AUC} & \textbf{Errors} \\
        \midrule
        initial-seed17 (frozen) & 17 & 0.0194 & 0.9992 & 1 \\
        initial-seed29          & 29 & 0.0143 & 0.9999 & 1 \\
        initial-seed43          & 43 & 0.0394 & 0.9998 & 3 \\
        lr3e-4-seed17           & 17 & 0.0261 & 1.0000 & 2 \\
        ep8-seed17              & 17 & 0.0203 & 0.9992 & 1 \\
        \bottomrule
    \end{tabular}
\end{table}

The frozen baseline's validation loss sits around the middle of the range
produced by changing the seed alone. That is the evidence missing from every
legacy result: this score is an ordinary outcome of a stable recipe, not a
number that happened to come out well once. One caution, which
Chapter~\ref{ch:multimodal} makes concrete: agreement between seeds
\emph{in-domain} says little about agreement \emph{across} corpora. On DFDC the
same three seeds turn out to be far apart.

\section{Measuring Generalization Instead of Assuming It}

Next we evaluated the frozen checkpoint, unchanged, on Celeb-DF-v2, a dataset
it never saw during training or threshold selection, and on a small locked
external benchmark, MNW~\cite{mnw}. Table~\ref{tab:crossdataset} puts these
results next to the in-distribution reference score.

\begin{table}[ht]
    \centering
    \caption{Current branch: the frozen baseline evaluated outside its training distribution, something no legacy stream was ever tested against.}
    \label{tab:crossdataset}
    \begin{tabular}{llrrr}
        \toprule
        \textbf{Dataset} & \textbf{Scope} & \textbf{Rows} & \textbf{ROC AUC} & \textbf{Balanced acc.} \\
        \midrule
        FakeAVCeleb (reference)     & development\_validation & 400 & 0.9992 & 0.9975 \\
        Celeb-DF-v2                 & generalization\_celebdf & 462 & 0.6476 & 0.5741 \\
        External benchmark (pilot)  & external\_mnw           & 85  & 0.0238 & 0.0952 \\
        External benchmark (base.)  & external\_mnw           & 85  & 0.0357 & 0.0952 \\
        \bottomrule
    \end{tabular}
\end{table}

On Celeb-DF-v2, ROC AUC drops from 0.999 to 0.6476, and balanced accuracy
drops from 0.9975 to 0.5741, which is close to chance. This is a real
measurement, not a hunch. It is also the first time in the project's history,
on either branch, that a trained checkpoint was scored on a dataset it was not
trained on. The external benchmark's score is even lower. The project's own
reporting flags it as most likely a label or evaluation mismatch, not a
confirmed result, and keeps it out of any firm interpretation until it has been
looked into. That kind of self-check was largely missing when the legacy branch
reported its perfect transformer score.

\section{Provenance: Making the DINOv2/DINOv3 Problem Structurally Impossible}

Before any checkpoint is allowed to produce a verdict, the current branch
checks six values: the checkpoint's SHA-256 hash, its MLflow run ID, the hash
of the data split used to validate it, the git commit of the training code, the
random seed, and a hash of the preprocessing configuration. If any of the six
differs from what was recorded at training time, the checkpoint refuses to
load. This was not a patch for one legacy bug. It is a general fix for the
whole class of problems that the DINOv2/DINOv3 drift belongs to, where a
document or folder name says something about a model that the model itself
does not confirm. Under the current system that kind of drift cannot survive
quietly, because nothing downstream will accept a checkpoint whose provenance
fails.

The same idea extends to cached inputs. The cache of preprocessed views is
keyed on a hash of the preprocessing settings and code version, and a model
will refuse to run on a cache that was built differently. When we later
retrained the shipped visual recipe from scratch, it reproduced the original
checkpoint exactly, epoch for epoch, which is good evidence that the whole
chain is deterministic.

\section{Reporting Discipline}

The current project also makes explicit distinctions that the legacy branch's
vocabulary did not have. A \emph{development result} is measured on a held-out
split of the training dataset. A \emph{generalization result} is measured on a
dataset never touched during training or tuning. We also separate
\emph{recorded} evidence from \emph{reproduced} evidence, depending on whether a
number in a report has been rerun in the current checkout or only exists from a
past run. The legacy evaluation log came close to this when it called the
transformer's perfect score ``an in-distribution ceiling, not evidence of
generalization.'' But that insight lived in a single closing note in one log
file. It was not part of a system that applied the distinction everywhere a
number appears.

\section{Dataset Roles on the Current Phase}

The legacy phase spread its attention across many manipulation categories
within one dataset. The current phase instead gives each dataset its own job,
and those jobs are not interchangeable. FakeAVCeleb is still the main
development dataset, used for training and for the seed comparison in
Table~\ref{tab:seed-comparison}. Celeb-DF-v2 is used only as a cross-dataset
generalization set; it is never touched during training or threshold selection
and exists to produce the measurement in Table~\ref{tab:crossdataset}. The
small locked external benchmark is reserved in the same way, for evaluation
only. When this baseline was frozen, FaceForensics++ was declared for future
visual experiments but had no videos downloaded in the working checkout. It was
later used as a second training corpus, as described in
Chapter~\ref{ch:multimodal}, which also adds DFDC as the main cross-corpus test.
Separating these roles, with one dataset to train on and others reserved for
testing generalization, is itself a structural improvement. The legacy phase
had correctly picked Deepfake-Eval-2024 for its final generalization test but
never actually ran it.

\section{The Dashboard}

The current dashboard deliberately does less than the legacy multi-page
dashboard, and that is a design choice, not a step backwards. Its Evidence gate
is the only page that produces a verdict, and it loads the frozen
EfficientNet-B0 checkpoint only after all six provenance checks pass. When a
clip is uploaded, the dashboard computes that clip's SHA-256 hash, and any new
upload clears whatever preprocessing or prediction was cached for the previous
clip, so a stale result can never be shown against the wrong input. If no
usable face sequence can be extracted, the dashboard returns an explicit
\texttt{indeterminate} verdict along with the reason, instead of forcing a real
or fake label out of insufficient evidence. Coverage is reported next to every
metric for the same reason. This follows directly from the reporting
discipline described above: a system that carefully separates kinds of
evidence in its write-up should not invent a confident verdict for a clip it
cannot actually read.

\section{Reproducible Environment}

The current phase's main checkout runs a CUDA-specific environment,
\texttt{torch 2.12.1+cu130}, verified against a detected GPU, and it can be
rebuilt with a single command:

\begin{verbatim}
uv sync --extra cu130 --extra ml --extra media --extra dashboard \
  --extra tracking --group dev
\end{verbatim}

The \texttt{cpu} and \texttt{cu130} extras are declared as conflicting and must
never be installed together. Doing so silently produces a Torch build that
reports the GPU as unavailable even on working hardware, which is exactly the
kind of quiet failure the provenance checks elsewhere are meant to prevent. The
legacy phase had a similar but narrower note (use exactly one of
\texttt{--extra cu130} or \texttt{--extra cpu}, never both), but its
documentation has no single-command rebuild tied to a specific, verified driver
and toolkit version.

% --
\chapter{Returning to the Multimodal Plan}
\label{ch:multimodal}

\section{Why We Went Back}

Once the visual baseline had a trustworthy number, the obvious next question
was the one the legacy phase set out to answer in the first place. Does adding
audio and cross-modal streams actually help? The argument for it is
convincing. Lip-sync forgeries change only the mouth and the voice, while the
background, the lighting and the upper face stay genuine, so a detector reading
pixels alone mostly sees real video. Reading both modalities and looking for
disagreement, such as a mouth forming one sound while the audio carries
another, should catch what a visual model misses.

We fixed the research question before running any of these experiments and
allowed for a negative answer: under source-disjoint, shortcut-controlled
evaluation, does cue-specific fusion of visual artifacts, audio spoofing and
mouth-audio alignment generalize better than a strong visual baseline? On our
data it does not. This chapter explains how we reached that answer, and the
four results that would have slipped past us if we had judged the system by
accuracy tables alone.

\section{Two Systems, One Evaluation Path}

\textbf{Design A} keeps the current visual branch and adds two more. The
visual branch is the EfficientNet-B0 and GRU model from
Chapter~\ref{ch:current}, reading 16 face crops. The audio branch uses
Wav2Vec2~\cite{wav2vec2} with attentive pooling over 4~s of 16~kHz audio to
detect spoofing. The sync branch pairs a framewise ResNet-18 mouth encoder with
a second Wav2Vec2 and reads a 2~s window of 50 mouth crops to judge whether the
mouth matches the sound. These are combined by a regularised logistic
regression over Platt-calibrated branch logits plus three quality features, and
that fusion model is fitted only on out-of-fold rows so it never sees scores
leaked from training.

\textbf{Design B} is much closer to the legacy plan. It replaces the branches
with configurable streams, and replaces the late fusion head with a
feature-level one. A visual stream is a backbone, either EfficientNet-B0 or a
frozen DINOv3 ViT-S/16~\cite{dinov3}, followed by a temporal model. An
audio-visual stream applies cross-attention between temporal tokens:
\begin{equation}
A = \operatorname{softmax}\!\left(\frac{(X_a W_Q)(X_v W_K)^\top}{\sqrt{d}}\right),
\qquad \tilde{X} = A\,(X_v W_V),
\end{equation}
where the audio tokens $X_a$ are the queries and the video tokens $X_v$ supply
the keys and values. The stream then embeds the residual $X_a - \tilde{X}$,
averaged over time. Put simply, this residual is whatever part of the sound the
mouth does not account for. The lip-sync stream reads the mouth window with its
audio, and the emotion stream reads the face crops with the voice. The fusion
head projects each stream's 256-dimensional embedding to a common width and
masks any missing stream \emph{after} the projection, so an absent stream
contributes exactly zero. It then classifies the concatenated vector with a
multi-layer perceptron. This is the design shown in
Figures~\ref{fig:stage2b} and~\ref{fig:stage3}, now built and trained.

Both designs are scored by the same evaluation code. When the largest gap in
our results first appeared, we suspected an evaluation bug. To check, we
scored the Design A checkpoint through the Design B pipeline, and it reproduced
the same number to four decimal places.

\section{A Check We Record but Never Optimize}

If a stream really aligns mouth movement with sound, its attention map $A$
should put most of its weight near the diagonal, since speech reaches the
microphone at an almost constant delay after the lip movement that produced it.
So after every epoch we record the diagonal mass
\begin{equation}
m = \frac{1}{T_q}\sum_{i,j} A_{ij}\,\mathbb{1}\!\left[\,\left|\tfrac{T_k}{T_q}\,i - j\right| \le 1\right]
\end{equation}
and we deliberately keep it out of the loss. Once a measure is being optimized,
it stops telling you anything about the mechanism and becomes one more thing
for the model to fit. The value you would get by chance depends on the query
length: uniform attention gives 0.0592 at 50 steps and 0.3438 at 8 steps.

\section{Experimental Setup}

The first track trains on FakeAVCeleb~\cite{fakeavceleb}, split 70/15/15 by
source identity. The main cross-corpus test set is DFDC~\cite{dfdc}, which is
the only corpus with audio used here that is not built on VoxCeleb2.
Celeb-DF-v2~\cite{celebdf} and the evaluation-only MNW benchmark~\cite{mnw}
serve as further external checks. LAV-DF~\cite{lavdf} is used to test the
cross-modal streams under a stricter setup, and FaceForensics++~\cite{ffpp}
c23 serves as a second training corpus following the published cross-dataset
protocol.

The second track uses DF26~\cite{df26}, which has 271 real clips and 2,420
generated clips from seven generators, four commercial and three open source.
A frozen encoder with a logistic probe is trained on six generators and scored
on the seventh.

The protocol fixes three training seeds: 17, 29 and 43. Every Design A branch
and the fusion head were trained with all three. Each Design B stream, however,
has only seed 17, and seed 17 turns out to be the best of the three on DFDC. A
single-seed cross-corpus number in this chapter should therefore be read as a
best case, not a typical one.

We report 95 percent bootstrap confidence intervals~\cite{bootstrap},
clustered on source identity, because what a reader really wants to know is
how the detector behaves on a different set of speakers. A resample that ends
up containing only one class is skipped and counted, never scored as 0.5. When
two systems are compared, the comparison is paired: both are scored on the
same resampled identities and the difference is taken within each resample.
For DF26 we resample clips or source videos, since generated clips have no
speaker.

A note on Celeb-DF numbers: the Celeb-DF figures in this chapter come from the
later evaluation pipeline and are not the same run as
Table~\ref{tab:crossdataset}. They should not be compared with it to the
fourth decimal place.

\section{The Cross-Modal Streams Do Not Read Correspondence}

On held-out FakeAVCeleb clips the emotion stream reaches 0.9991 ROC AUC, the
best of the five Design B streams. Table~\ref{tab:diag} shows the diagonal mass
for this stream and the lip-sync stream over every epoch of three training
runs.

\begin{table}[ht]
    \centering
    \caption{Diagonal attention mass against its uniform (chance) value.}
    \label{tab:diag}
    \begin{tabular}{lrrl}
        \toprule
        \textbf{Stream} & \textbf{Steps} & \textbf{Uniform} & \textbf{Measured} \\
        \midrule
        emotion & 8 & 0.3438 & 0.3434 to 0.3457, 9 epochs \\
        lip-sync & 50 & 0.0592 & 0.0592 to 0.0595, 4 epochs \\
        lip-sync, lower encoder LR & 50 & 0.0592 & 0.0592 to 0.0595, 6 epochs \\
        \bottomrule
    \end{tabular}
\end{table}

Over nine epochs the emotion stream's ROC AUC rose from 0.9847 to 0.9991, while
its attention barely moved from where it started. The accuracy is real, but the
mechanism the stream is named after is not doing the work.

The likely explanation lies in the data. FakeAVCeleb's fakes are manipulated
across the whole clip, so a model can separate them by recognising the
generator without ever comparing mouth and sound. LAV-DF lets us block that
route. From each file we cut two windows, one containing the manipulated
segment and one avoiding it, so speaker, lighting, microphone and codec are the
same within the pair. A lip-sync stream trained on these pairs reaches 0.9969
ROC AUC on 1,161 held-out windows, with a diagonal mass of 0.0592 against a
uniform value of 0.0600. Matching the windows removes file-level cues, but
local artifacts from the manipulation can still separate the classes.
Near-uniform diagonal mass means we cannot claim the model learned temporal
correspondence.

\section{The Visual Models Partly Read the Video Codec}

In FakeAVCeleb, genuine videos are H.264 at about 150~kbps. Every clip in the
face-swap plus Wav2Lip category, the largest fake category, is MPEG-4 Part~2 at
about 330~kbps. We re-encoded 75 test clips for each row and scored them with
the deployed ensemble at a threshold of 0.5. Table~\ref{tab:codec} shows what
happened.

\begin{table}[ht]
    \centering
    \caption{Share of clips called fake after re-encoding (75 clips per row).}
    \label{tab:codec}
    \begin{tabular}{llr}
        \toprule
        \textbf{Clips} & \textbf{Codec} & \textbf{Called fake} \\
        \midrule
        genuine & original H.264 & 0.0\% \\
        genuine & re-encoded H.264 (control) & 0.0\% \\
        genuine & re-encoded MPEG-4 & \textbf{93.3\%} \\
        swap plus Wav2Lip & original MPEG-4 & 100.0\% \\
        swap plus Wav2Lip & re-encoded H.264 & \textbf{30.7\%} \\
        face swap, H.264 & re-encoded MPEG-4 & 100.0\% \\
        \bottomrule
    \end{tabular}
\end{table}

The control row is what makes the table readable. Re-encoding a genuine clip
with its own codec changes its pixels but does not change the verdict at all.
Changing only the codec flips the verdict in both directions. Since decoding
returns exactly the pixels the encoder wrote, compression artifacts included, a
model can end up learning the encoder. This also casts doubt on the legacy
phase's in-distribution scores in Table~\ref{tab:legacy-results}: those
streams were trained and tested on the same FakeAVCeleb codec split, so some of
their accuracy may come from this cue as well, although we have not run the
codec-swap test on the legacy checkpoints.

\section{Training Against the Codec Cue Costs Transfer}

To find out what happens when the model is pushed away from the codec cue,
we trained four recipes of the EfficientNet-B0 model at seed 17, each changing one thing relative to
the one before. The first is the shipped recipe, selected on ROC AUC. The
second adds random JPEG, H.264 and MPEG-4 round trips to both classes and
selects checkpoints on codec-swapped validation clips. The third adds
self-blended fakes~\cite{sbi}, and the fourth adds the FaceForensics++
training split.

\begin{table}[ht]
    \centering
    \caption{Visual recipes against the codec cue (ROC AUC).}
    \label{tab:recipes}
    \begin{tabular}{lrrr}
        \toprule
        \textbf{Recipe} & \textbf{FakeAVCeleb} & \textbf{DFDC} & \textbf{Celeb-DF} \\
        \midrule
        shipped & 0.9988 & \textbf{0.7611} & 0.6682 \\
        + compression & 0.9959 & 0.5720 & 0.6525 \\
        + self-blending & 0.9978 & 0.6496 & 0.6762 \\
        + FaceForensics++ & 0.9904 & 0.6756 & 0.6509 \\
        \bottomrule
    \end{tabular}
\end{table}

Every modified recipe did worse than the shipped one on DFDC
(Table~\ref{tab:recipes}). Compression hurt most, by 0.1891, with
non-overlapping intervals. Self-blending and the larger corpus won some of that
back but still sit 0.0855 below the shipped recipe. On Celeb-DF the intervals
overlap for all four, so none of them can be told apart. One caveat matters
here: compression augmentation and checkpoint selection changed at the same
time, so this experiment cannot isolate how much of the transfer loss comes
from removing the codec cue itself.

\section{Fusion Does Not Beat the Visual Branch Across Corpora}

Table~\ref{tab:fusion} gives the paired difference between the fused Design A
system and a head fitted on the visual branch plus the same quality features.
Each seed uses its own cross-fitted fold models.

\begin{table}[ht]
    \centering
    \caption{Fusion minus visual, paired, with 95\% intervals.}
    \label{tab:fusion}
    \begin{tabular}{lll}
        \toprule
        \textbf{Seed} & \textbf{FakeAVCeleb test} & \textbf{DFDC} \\
        \midrule
        17 & $+0.0159$ [$+0.0129$, $+0.0190$] & $-0.0085$ [$-0.0196$, $+0.0027$] \\
        29 & $+0.0170$ [$+0.0143$, $+0.0195$] & $-0.0215$ [$-0.0371$, $-0.0052$] \\
        43 & $+0.0184$ [$+0.0159$, $+0.0209$] & $-0.0417$ [$-0.0698$, $-0.0129$] \\
        \bottomrule
    \end{tabular}
\end{table}

In-domain, fusion helps at every seed. Almost all of that gain comes from 75
clips where the video is genuine and only the voice was replaced, which no
visual branch could ever catch. Across corpora, fusion never comes out ahead
and is measurably worse at two seeds.

Design B tells the same story. We fitted every subset of the five streams on
the same rows. The best combination on DFDC, emotion plus lip-sync, scores
0.6076, while fusing all five scores 0.5411, lower than three of its own
inputs. Across corpora the audio branch is at chance, 0.5054 [0.4747, 0.5371],
and the sync branch's anomaly score stays below 0.5 at all three seeds. In
short, one branch transfers and the other two do not.

\section{Seed Spread Exceeds the Measured Fusion Differences}

\begin{table}[ht]
    \centering
    \caption{DFDC ROC AUC by training seed.}
    \label{tab:seeds}
    \begin{tabular}{lrrrr}
        \toprule
        \textbf{System} & \textbf{17} & \textbf{29} & \textbf{43} & \textbf{Spread} \\
        \midrule
        visual branch & 0.7611 & 0.6963 & 0.6878 & 0.0733 \\
        fused system & 0.7500 & 0.6698 & 0.6422 & \textbf{0.1078} \\
        \bottomrule
    \end{tabular}
\end{table}

In-domain the seeds agree to three or four decimal places, as
Table~\ref{tab:seed-comparison} already showed. Across corpora the fused system
ranges over 0.1078 (Table~\ref{tab:seeds}), about thirteen times the seed-17
fusion difference. Fusing three branches made the seed noise worse instead of
averaging it out.

Averaging helped when it was done across checkpoints instead of branches. An
equal-weight average of the six final visual checkpoints (EfficientNet-B0 and
ConvNeXt-Tiny~\cite{convnext}, each at three seeds) scores 0.7739
[0.7401, 0.8064] on DFDC. That is $+0.0469$ [$+0.0286$, $+0.0666$] over the
EfficientNet seed mean and level with its best seed; on Celeb-DF it is level
with every member. We should be clear that this grouping was chosen after
comparing three groupings on DFDC, so its DFDC gain still has to be confirmed
on a test partition we have not looked at.

\section{Choices That Looked Right on Validation}

\textbf{Selecting on loss.} One visual stream reached its lowest validation
loss at epoch 1, while its backbone was still frozen. That checkpoint scored
0.4668 ROC AUC in-domain, which is below chance. Binary cross-entropy punishes
confident mistakes heavily, so a model that ranks well but is overconfident can
lose to one that hedges and ranks badly. Since then every trainer in the
project selects on validation ROC AUC.

\textbf{Frozen BatchNorm.} Freezing BatchNorm statistics during fine-tuning
raised validation from 0.5154 to 0.9058, and we first logged it as a clear
win. Compared with frozen statistics, live statistics lost 0.5430 of
in-domain ROC AUC but gained 0.1276 on DFDC. However, the two arms had run for
different numbers of epochs. With the training budget held fixed, the DFDC
difference shrinks to 0.0177, about a seventh of the uncontrolled figure. Live
statistics act as an accidental domain adapter~\cite{adabn}, and the
in-domain cost is the part of the effect that survives the control. This is
also the clearest case of our own checks catching our own mistake: we had
already written the larger number down.

\section{The Operating Point Fails Before the Ranking Does}

\begin{table}[ht]
    \centering
    \caption{Visual branch under corruption of the cached view.}
    \label{tab:corruption}
    \begin{tabular}{lrrr}
        \toprule
        \textbf{Condition} & \textbf{ROC AUC} & \textbf{Genuine called fake} & \textbf{Fakes missed} \\
        \midrule
        clean & 0.9988 & 0 of 150 & 3 of 1,498 \\
        JPEG 75 & 0.9934 & \textbf{76 of 150} & 1 \\
        JPEG 50 & 0.9187 & \textbf{150 of 150} & 0 \\
        \bottomrule
    \end{tabular}
\end{table}

JPEG compression at quality 75 costs only 0.0054 of ROC AUC, which in a table
looks like robustness. In practice it pushes false alarms from 0 to 51 percent
(Table~\ref{tab:corruption}). Audio noise at 20~dB SNR breaks things the other
way: ROC AUC drops by 0.0227, but 857 of 859 fakes are missed. Across corpora
the model is also more confident when it is wrong than when it is right. On
DFDC the visual branch gets 1,942 of 2,410 clips wrong at a 0.5 threshold, with
a mean confidence of 0.9244 on its mistakes and 0.8517 on its correct answers.
No threshold on that confidence removes errors faster than it removes correct
answers. At 0.5, the visual branch calls 141 of 155 genuine Celeb-DF videos
fake. This is the legacy log's warning about saturated predictions, now
measured on data the model has never seen.

One partial fix works well. We fitted a Gaussian to frozen DINOv3 features of
genuine training clips and score each new clip by its Mahalanobis distance. This
model never sees a fake; it only asks whether the footage looks familiar. At
the deployed setting it abstains on 8.7 percent of held-out genuine
FakeAVCeleb clips and intercepts 74.4 percent of genuine Celeb-DF clips, with a
range of 49.4 to 83.1 percent over twenty reference splits. That brings false
accusations on unfamiliar footage down from about 91 percent to somewhere
between a quarter and a third. It fits naturally with the dashboard's
\texttt{indeterminate} verdict: when the system does not recognise the
footage, it should say so instead of guessing.

\section{Open-World Generated Video}

The second track applies the same discipline to video generated whole. Before
reporting any DF26 number we found and removed four shortcuts: scene content,
real clips shared across splits, sibling modes of the same generator, and
audio presence. The last one is enough on its own to predict the label, because
every commercial clip carries an audio track and almost no other clip does. Real
clips were also the only ones at 25 or 29.97 frames per second, so we
re-encoded every clip to $640\times360$ at 24 frames per second.

\begin{table}[ht]
    \centering
    \caption{Unseen-generator ROC AUC on DF26, frozen encoder and linear probe.}
    \label{tab:df26}
    \begin{tabular}{lrl}
        \toprule
        \textbf{Encoder} & \textbf{Macro} & \textbf{Paired against DINOv3} \\
        \midrule
        Perception Encoder~\cite{pe} & \textbf{0.7527} & $+0.0741$ [$+0.0121$, $+0.1317$] \\
        CLIP ViT-B/16~\cite{clip} & 0.7454 & $+0.0668$ [$+0.0049$, $+0.1289$] \\
        DINOv3 ViT-S/16 & 0.6786 & reference \\
        \bottomrule
    \end{tabular}
\end{table}

Perception Encoder leads on all seven held-out generators
(Table~\ref{tab:df26}), from 0.8496 on Veo~3.1 down to 0.6542 on LTX~2.3. There
is an important caveat. All three encoders were given ImageNet-normalized
pixels, which is right for DINOv3 but wrong for the other two. These numbers
were therefore measured under mismatched preprocessing and should not be
treated as lower bounds. The comparison stays provisional until the features
are re-extracted with each encoder's own normalization.

Which design works depends on which part of the frame the model looks at. The
same full-frame probe on FaceForensics++, with one manipulation family held out
in each run, is at chance, with a macro score of 0.5387. Going the other way,
face crops through the fine-tuned visual stream, retrained with one family
held out each time, lose 0.1231 macro against a model that saw all six, and the
graphics-based FaceSwap method falls to 0.5195. Full frames suit video that is
generated whole; face crops suit a real face that has been manipulated.

\section{The Encoders the Legacy Plan Named}

The legacy plan specified AV-HuBERT~\cite{avhubert} for lip-sync and
HSEmotion~\cite{hsemotion} for the face side of the emotion stream, and neither
was ever integrated on the old branch. Neither runs in our current environment
as released, either. AV-HuBERT depends on fairseq, which pins an older version
of PyTorch, and HSEmotion ships as a pickled model that current timm cannot
load. So we rewrote AV-HuBERT Base in plain PyTorch and loaded HSEmotion's
weights into the same EfficientNet-B0 architecture. We checked both against the
official code on identical inputs, and their features and logits agree to
within $2\times10^{-6}$. AV-HuBERT is run twice, once on the lips alone and
once on the audio alone, because its pretraining placed both modalities in the
same space.

On 98 percent of 128 genuine training clips, correctly aligned audio matches
the lips better than audio shifted by half a second. A lip-sync stream built on
these features reaches 0.9416 validation ROC AUC, compared with 0.7593 for the
Wav2Vec2-based stream. The selected checkpoint is epoch 2 of five, before the
encoder is unfrozen at epoch 3. Its diagonal mass stays between 0.05875 and
0.05919, very close to the uniform value of 0.0592. So the validation gain gives
no evidence of diagonal alignment in the trained attention head. The
shifted-audio check says something about the pretrained features, not about the
mechanism of the detector built on top of them.

These results are preliminary and use seed 17 only. Held-out FakeAVCeleb, DFDC
and fusion scores for the replacement encoders are still pending, and the
HSEmotion training run has not produced a completed result for this
comparison. Validation accuracy alone cannot show that either replacement
improves cross-corpus detection.

\section{What This Chapter Adds}

Four results here would each have gone unnoticed if we had looked only at
accuracy tables. A cross-attention stream reached near-perfect accuracy while
its attention never left its starting value. The visual models partly learned
the codec. A selection rule kept a checkpoint that scored below chance. And the
training seed moved the cross-corpus result more than any architecture choice
we tested. Each was caught by a quantity recorded next to the objective and
never optimized against it. The legacy branch had no such quantities, which is
why none of these questions could even be asked there.

% --
\chapter{Legacy vs. Current: A Direct Comparison}
\label{ch:comparison}

\section{Side by Side}

Table~\ref{tab:headtohead} compares the two branches on the things that decide
whether a detector's reported performance can be trusted, not just on how high
the headline number is.

\begin{table}[ht]
    \centering
    \caption{Legacy branch versus current branch, dimension by dimension.}
    \label{tab:headtohead}
    \begin{tabular}{p{3.6cm}p{5.4cm}p{5.4cm}}
        \toprule
        \textbf{Dimension} & \textbf{Legacy (\texttt{old} branch)} & \textbf{Current (\texttt{main} branch)} \\
        \midrule
        Visual streams trained & Three: EfficientNet-B0, Xception, self-supervised transformer & EfficientNet-B0 + GRU baseline repeated five times; later EfficientNet-B0, ConvNeXt-Tiny and frozen DINOv3 streams \\
        Best in-distribution ROC AUC & 1.000 (flagged by its own log as an unverified ceiling) & 0.9992 (stable across four additional seed/hyperparameter runs) \\
        Seed repetition & None, every stream trained exactly once & Baseline and every Design A branch at seeds 17, 29 and 43; Design B streams at seed 17 only \\
        Cross-dataset evaluation & None ever run, though a held-out dataset was named for this purpose & Celeb-DF-v2, DFDC, MNW, plus LAV-DF and DF26 for specific tests \\
        Cross-dataset result & Unknown & 0.6476 ROC AUC on Celeb-DF-v2 for the frozen baseline; 0.7611 to 0.6878 on DFDC depending on seed \\
        Audio-visual streams & Lip-sync and emotion streams: empty folders, never coded & Built and trained; high in-domain accuracy, but attention stays at the uniform value, so the cross-modal mechanism is unverified \\
        Fusion & Empty folder, never coded & Trained at three seeds; better in-domain, never better than the visual branch on DFDC \\
        Shortcut checks & None & Codec-swap test, silence offset, DF26 audio-presence and frame-rate controls \\
        Provenance checking & None; checkpoints loaded by path alone & Six-value check (hash, run ID, split hash, commit, seed, preprocessing hash) before any load \\
        Experiment tracking & Weights and Biases, aggregate and per-category metrics logged, no failed-run retention policy stated & MLflow, including failed runs kept deliberately for their debugging history \\
        Naming consistency & Checkpoint labeled \texttt{dinov2}, documentation and folder name say ``DINOv3'' & Every number traceable to one exact MLflow run ID and checkpoint hash \\
        Reporting vocabulary & No formal distinction between in-distribution and generalization claims & Explicit development-result / generalization-result / recorded-vs-reproduced-evidence distinctions used throughout \\
        \bottomrule
    \end{tabular}
\end{table}

\section{So Is the New Model Actually Better?}

The answer has two parts, and both matter.

On in-distribution accuracy alone, the two branches are too close for either to
claim a clear win. The legacy EfficientNet-B0 stream reached 0.996 test ROC
AUC, and the current EfficientNet-B0 plus GRU baseline reaches 0.9992
validation ROC AUC. These came from different splits, class ratios and
evaluation protocols, so calling the half-point gap an improvement would be
exactly the kind of overclaiming this report argues against. The legacy
transformer's 1.000 AUC is higher than anything on the current branch, but its
own log already warned against trusting it, and nobody checked whether that
warning was justified. The codec experiment in Chapter~\ref{ch:multimodal}
gives one more reason for caution: on FakeAVCeleb, near-perfect in-domain scores
can come partly from the video encoder, not the manipulation.

On the question that actually matters, whether a model can be trusted on data
it was not trained on, the current branch is clearly ahead, for a simple
reason. It is the only branch that has ever measured this. The legacy streams
might generalize well, or as poorly as the current baseline's 0.6476 on
Celeb-DF-v2, or worse; there is no way to know, because the measurement was
never taken. The current model is not necessarily better in some abstract
sense. What has changed is that its real strengths and weaknesses are now
known, written down and checked automatically every time the checkpoint loads.

The multimodal experiments sharpen this point. If the legacy plan had been
built as written and judged the way the legacy branch judged its visual
streams, it would most likely have looked like a success. The emotion stream
alone scores 0.9991 in-domain, and fusion adds $+0.0159$ at every seed. Only the
checks the current phase added, cross-corpus testing, seed spread and the
diagonal attention mass, show that the fused system does not beat the visual
branch on new data and that its cross-modal streams do not compare mouth and
sound the way their names suggest. That is the improvement this report
documents, and it is worth more than a slightly higher accuracy score.

\section{What This Means Going Forward}

The comparison does not mean the legacy ambitions should be dropped. The
five-stream, cross-attention, fully fused architecture is still a reasonable
goal, and the current phase has now built most of it. What it means is that the
current branch's discipline has to come along with it, instead of repeating the
old pattern of building outward faster than any one piece can be checked. In
practice, any future visual backbone, whether that means revisiting the legacy
Xception and DINOv2 checkpoints or training new ones, should go through seed
repetition, cross-dataset testing and the codec-swap test before its
in-distribution score is believed. Any future cross-modal stream should report
its diagonal mass next to its accuracy, and should be tied into the six-value
provenance check from its first checkpoint. The legacy branch showed the
five-stream idea could be built. The current branch showed that a trained
system's number means very little until you know how it behaves outside its own
training data, and whether its parts do what they are named for. The next phase
of this work should combine both lessons instead of choosing between them.

% --
\chapter{Current Verification State}
\label{ch:verification}

\section{Introduction}

A report that spends this much time on the gap between what is claimed and
what is true needs to apply the same standard to itself. This chapter sets out
what is currently verified about the current phase, what is blocked, and what
is still pending, so that the earlier chapters do not suggest a level of
readiness the codebase has not reached.

\section{What Is Currently Verified}

The CUDA environment installs cleanly and reports the GPU as available on the
detected hardware. The provenance-checked loader loads the frozen visual
checkpoint with all six checks passing: checkpoint hash, MLflow run ID, split
hash, training commit, seed and preprocessing hash. The Evidence gate runs end
to end. Given a synthetic clip with no detectable face, it returns an
indeterminate verdict with a named blocker instead of guessing, which confirms
that the verdict logic follows actual evidence coverage. The automated test
suite, linting and documentation build all pass. Retraining the shipped visual
recipe reproduces the original checkpoint exactly, epoch for epoch, and the
Design A checkpoint scored through the Design B pipeline matches its original
number to four decimal places. The multimodal results in
Chapter~\ref{ch:multimodal} are registered with the SHA-256 hashes of their
artifacts.

\section{What Is Currently Blocked}

Two things cannot currently be reproduced from the active checkout, and we
want to say so directly. The first is a fresh inference pass on real video. The
raw dataset folder referenced by the project's manifest is empty in the working
checkout, so no real manifest row can currently be decoded end to end through
the live pipeline. A single-clip probability recorded earlier gives us a
specific target to reproduce once the dataset is restored, which is better than
a vague plan to ``verify things later.'' The second is that the cross-dataset
numbers in Table~\ref{tab:crossdataset} are recorded evidence from an earlier
run, not something we reran for this report, because the datasets and cache
needed for that exact evaluation are not currently in a state we could
re-execute ourselves.

\section{What Is Still Pending or Limited}

Several limits apply to the multimodal results in particular.

\begin{itemize}
    \item Every Design B stream has one training seed where the protocol asks
    for three. Table~\ref{tab:seeds} shows how much a single seed can be worth
    across corpora, and seed 17 is the most favourable of the three on DFDC.
    \item Design A and Design B both train on FakeAVCeleb, so their numbers
    compare with each other, not with published tables. Only the
    FaceForensics++ result follows the published protocol, and it too uses one
    seed.
    \item None of the training data comes from the wild.
    \item The cross-modal streams do not measure correspondence, so the
    mechanism the architecture is named after remains unverified.
    \item One exit criterion from the research design, that the sync offset
    interval should exclude chance, was never measured, because the sync
    branch's training history recorded only a combined loss.
    \item The first track is nearly blind to fully generated video: it detects
    0 of 10 MNW clips from one talking-head generator.
    \item The second track rests on one corpus and one split, its encoder
    comparison used mismatched normalization, and DF26 cannot support audio
    work.
    \item The six-checkpoint ensemble's DFDC gain was measured after choosing
    its grouping on DFDC, and the AV-HuBERT and HSEmotion results have no
    held-out evaluation yet.
\end{itemize}

\section{Why the Distinction Is Kept Visible}

It would be easy to blur the line between a number that has been freshly
reproduced and one that has only been recorded from a past run, since both end
up looking the same in a results table. We keep them apart because the whole
point of the provenance system in Chapter~\ref{ch:current} is that every number
should trace back to exactly what produced it. Recorded evidence still has that
traceability, since it points to a specific MLflow run ID and metrics file, but
it has not yet passed the stronger test of being regenerated in the current
session. Closing that gap, by restoring the dataset and reproducing the known
single-clip probability, is the first item in Chapter~\ref{ch:future}.

% --
\chapter{Future Scope}
\label{ch:future}

\section{Introduction}

The comparison in Chapter~\ref{ch:comparison} and the limits in
Chapter~\ref{ch:verification} point to a fairly specific list of next steps,
not an open-ended research agenda. We list them roughly in the order we think
they should be tackled.

\begin{itemize}
    \item \textbf{Restore the raw dataset and reproduce the known single-clip
    probability.} This comes first because it turns the current phase's
    cross-dataset evidence from recorded into reproduced, closing the gap
    described in Chapter~\ref{ch:verification}.
    \item \textbf{Train every Design B stream at seeds 29 and 43.} The protocol
    already calls for three seeds, and Table~\ref{tab:seeds} shows that a
    single cross-corpus number can sit anywhere in a 0.1078 range.
    \item \textbf{Confirm the six-checkpoint ensemble on an untouched test
    partition.} Its DFDC gain was measured after choosing the grouping on DFDC,
    so it needs a fresh test before we rely on it.
    \item \textbf{Separate codec augmentation from checkpoint selection.} The
    recipe comparison changed both at once. Running them as separate arms
    would show how much of the 0.1891 DFDC loss comes from removing the codec
    cue itself.
    \item \textbf{Revisit the legacy Xception and \texttt{dinov2} checkpoints
    under the current protocol.} Both reached strong in-distribution scores
    that were never checked for seed variance, cross-dataset performance or
    codec dependence. Running them through the same checks would finally show
    whether either legacy stream generalizes better, worse or about the same
    as the current baseline. The naming question is partly settled, since
    Design B now uses a real frozen DINOv3 ViT-S/16, but the legacy checkpoint
    should still be relabelled to match what it is.
    \item \textbf{Re-extract the DF26 features with each encoder's own
    normalization.} Until then, the Perception Encoder and CLIP results remain
    provisional.
    \item \textbf{Finish the held-out evaluation of the AV-HuBERT and HSEmotion
    streams.} That means held-out FakeAVCeleb, DFDC and fusion scores, plus the
    sync offset criterion that was never measured.
    \item \textbf{Look into the external benchmark's inverted-looking result.}
    An ROC AUC as far below 0.5 as the one recorded on the locked external
    benchmark usually points to a label or evaluation mismatch rather than
    genuinely bad performance, and it needs a targeted check before it is read
    either way.
    \item \textbf{Add training data from the wild and widen the second track.}
    Both tracks are limited by their data: track one has no in-the-wild
    training data, and track two rests on one corpus and one split.
    \item \textbf{Make seed repetition, cross-dataset testing and mechanism
    checks a standing requirement.} The most transferable lesson here is about
    procedure. No branch's in-distribution score should be treated as final
    until it has survived these checks, however promising it looks after the
    first training run.
\end{itemize}

\section{A Note on Priority}

The ordering is deliberate. The first few items cost little new engineering
and directly answer questions this report has already raised. The middle items
(the legacy checkpoints, DF26 normalization, the replacement encoders and the
external benchmark) are bounded pieces of work. The last two are a broader
return to the legacy phase's original ambition of a detector that works in the
real world. They come last because they should only start once the discipline
described in Chapters~\ref{ch:current}, \ref{ch:multimodal}
and~\ref{ch:verification} is simply how this project works, not an extra step
added at the end.

% --
\chapter{Conclusion}
\label{ch:conclusions}

We set out to compare the project's old trained models with its new ones, and
the comparison turned out not to be about which one scores higher. The legacy
branch's three visual streams were genuinely trained and evaluated, and their
in-distribution scores are, if anything, slightly higher than the current
baseline's. Its most extreme result, a perfect 1.000 AUC, was flagged by its
own authors as unproven instead of being presented as an achievement, which
shows real self-awareness in an unfinished project. What the legacy branch
never did, across three trained streams and months of plans for two more, was
find out whether any of its numbers meant anything outside the one dataset they
were all measured on.

The current branch did that. For the visual baseline the answer was much less
flattering than 0.999 but far more useful: a 35-point drop in ROC AUC as soon
as the evaluation dataset changed. That drop would very likely have applied to
the legacy streams too, and it fits what the literature in
Chapter~\ref{ch:literature} predicts for detectors tested outside their
training distribution \cite{ramanaharan2025, coccomini2022}. When we then built
the multimodal system the legacy branch had only planned, the same checks
showed that cue-specific fusion does not beat the visual baseline across
corpora on our data, and that fusing every stream was worse than three of its
parts. They also showed a cross-attention stream scoring 0.9991 with a
mechanism that does nothing, a visual model reading the codec, a regularizer
that bought validation accuracy by switching off a domain adapter, and a
selection rule that preferred a below-chance checkpoint. Every one of these was
caught by a quantity recorded next to the objective and never optimized
against it. The same habit caught our own mistakes, including an effect we had
already reported that shrank sevenfold once its control ran.

Reporting a single cross-corpus number to four decimal places invites
comparisons the measurement cannot support when the seed alone moves it by
0.1078. Together with a provenance system built so that a checkpoint's identity
can never again drift from what a document claims about it, the way the
legacy DINOv2 checkpoint drifted from its ``DINOv3'' documentation, this makes
the current branch something different from a bigger or flashier version of the
old one. It is the first version of the project that actually knows what it has
built.

% --
\addcontentsline{toc}{chapter}{Bibliography}
\bibliographystyle{IEEEtran}
\bibliography{references}

\end{document}